% 【本文件写什么】一级组件 wateros-abi（§ 级聚合）
%   - 组件职责与在内核中的位置：用户态与内核 ABI 约定：系统调用号、寄存器/栈约定、错误码映射。
%   - 聚合 crate os/components/wateros-abi/src/lib.rs 的 pub mod 树与对外导出面
%   - 子模块划分、与相邻组件的依赖边界（谁依赖谁、走哪条门面）
%   - 根 feature / 本组件 Cargo.toml feature 如何选用各 impl
%   - 1～2 段综述后，由下方 \input 展开子模块；不在此逐条列举 api trait
% 【事实来源】
%   - os/components/wateros-abi/src/lib.rs、Cargo.toml
%   - docs/exports/public-api/wateros-abi.md
%   - docs/exports/features/wateros-abi.md
%   - os/feature-tree.txt 中 wateros-abi 相关段
% 【不写什么】
%   - api-v0 契约细节 → 子目录 api/api-v0.tex
%   - 各 impl 算法/平台细节 → 子目录 impl/*.tex

\section{wateros-abi}

\code{wateros-abi}定义用户态与内核共享的系统调用ABI：Linux风格错误码、六个寄存器参数包、调用号抽象与\code{-errno}返回值编码。聚合\code{src/lib.rs}按feature挂载\code{api-v0}契约与具体号表实现；\code{default} feature为空，调用方须显式启用\code{api-v0}及\code{impl-linux-generic64}（或占位\code{impl-dummy}）。

主线平台feature经根内核传递\code{abi/impl-linux-generic64}，同时启用\code{api-v0}与\code{LinuxGeneric64}号表，并导出\code{ActiveSyscallNumberTable}别名。\code{wateros-syscall}、\code{wateros-platform} trap路径通过本组件读取\code{SyscallArgs}与符号化调用号；号表存在某常量不等于内核已实现对应\code{sys\_*}。

% 【本文件写什么】独立子 crate（无 api/impl 目录树）：abi-api
%   - 职责：ABI 顶层 api-v0 契约 crate。
%   - 源码：os/components/wateros-abi/abi-api/src/
%   - 对外暴露的函数/类型、在聚合 lib.rs 中的重导出方式
%   - 实现要点、feature（若有）、与其它子 crate 的边界
% 【事实来源】
%   - os/components/wateros-abi/abi-api/、父组件 src/lib.rs

\subsection{abi-api}

子crate\code{abi-api/api-v0}（包名\code{wateros-abi-api-v0}）提供四组模块：\code{errno}、\code{user\_ret}、\code{syscall\_number}、\code{syscall\_args}。聚合层在\code{api-v0} feature下按同名子模块再导出，便于\code{use abi::errno::ErrNo}等统一引用。

本crate为纯契约，\code{\#![no\_std]}，依赖\code{wateros-base-config}的\code{MAX\_SYSCALL\_ARGS}决定参数包槽位数。不包含任何平台号表具体数值；号表由\code{abi-impl}提供。

% 【本文件写什么】api-v0 契约层，叶子，只写语义不写实现
% - crate 路径：os/components/wateros-abi/abi-api/api-v0/src/
% - 列出并说明：pub trait、核心类型、错误枚举、常量
% - 每个公开 API 的前置条件、后置条件、线程/中断上下文限制
% - 与 Linux/用户态约定的对齐点与 deliberate 差异
% - v0 版本当前行为与后续可替换 extension 点
% 【事实来源】
% - 上述 crate 下全部 .rs 源文件
% - docs/exports/public-api/ 中对应符号表，以源码为准
% 【不写什么】
% - impl 内部数据结构、汇编、具体算法步骤

\paragraph{ErrNo}

\code{ErrNo(isize)}封装Linux errno正数值，与\code{asm-generic}数值一致。\code{raw()}返回正数errno；\code{user\_ret()}返回\code{-errno}供用户态\code{isize}承载。关联常量覆盖\code{EPERM}、\code{ENOENT}、\code{EINTR}、\code{EFAULT}、\code{EINVAL}、\code{ENOSYS}、\code{EAGAIN}、\code{ENOMEM}、\code{EACCES}、\code{EBADF}、\code{ECHILD}、\code{ENOSPC}、\code{ENOTTY}、\code{ENOTDIR}、\code{EISDIR}、\code{ENAMETOOLONG}、\code{ELOOP}、\code{ERANGE}等常用子集。\code{KernelResult<T>}为\code{Result<T, ErrNo>}别名。

\paragraph{UserRet 与 SyscallResult}

\code{UserRet(isize)}表示用户态可见的单一返回值：非负为成功，负值为\code{-errno}。\code{from\_success(usize)}、\code{from\_error(ErrNo)}、\code{from\_kernel\_result(SyscallResult)}完成编码。\code{SyscallResult = Result<usize, ErrNo>}是内核handler常用中间类型，由syscall层最终转为\code{UserRet}。

\paragraph{SyscallNumber 与 SyscallNumberTable}

\code{SyscallNumber(usize)}是调用号newtype，\code{new}/\code{raw}不校验内核是否实现。\code{SyscallNumberTable} trait以关联常量形式声明符号化调用号，按类别分组：文件与描述符、路径与xattr、进程与执行、内存、时间、凭证、信号、socket、多路复用等。trait只表达编号映射，不表示handler已实现；早期维护busybox/simple进程所需子集，随strace缺口逐步补齐。

\paragraph{SyscallArgs 与 SyscallPacket}

\code{SyscallArgs}为\code{repr(C)}的\code{[usize; MAX\_SYSCALL\_ARGS]}包装，当前 \code{MAX\_SYSCALL\_ARGS = 6}，顺序与ABI参数寄存器一致。\code{from\_regs}、\code{arg(idx)}、\code{as\_regs}提供构造与读写字段。\code{SyscallPacket}组合\code{SyscallNumber}与\code{SyscallArgs}，\code{new}不校验语义一致性。索引越界或未定义槽位语义属于调用方错误，由具体handler决定返回\code{EINVAL}或\code{EFAULT}。

\paragraph{与 Linux 对齐点}

错误码数值、成功非负/失败\code{-errno}约定、六个参数寄存器上限与Linux generic 64位用户态ABI对齐。deliberate差异：\code{SyscallNumberTable}未覆盖全部Linux调用；\code{SELECT}在generic64表上无独立nr，见 \code{impl-linux-generic64}。启用\code{impl-linux-generic64}时聚合层额外导出\code{ActiveSyscallNumberTable}类型别名。


% 【本文件写什么】独立子 crate，无 api/impl 目录树：abi-impl
% - 职责：ABI 实现容器：dummy 占位与 linux-generic64 主线。
% - 源码：os/components/wateros-abi/abi-impl/src/
% - 对外暴露的函数/类型、在聚合 lib.rs 中的重导出方式
% - 实现要点、feature，若有、与其它子 crate 的边界
% 【事实来源】
% - os/components/wateros-abi/abi-impl/、父组件 src/lib.rs

\subsection{abi-impl}

\code{abi-impl}容器承载两类实现：\code{impl-dummy}为工作区依赖占位；\code{impl-linux-generic64}为主线Linux\code{asm-generic}风格64位调用号表。根\code{Cargo.toml}的\code{impl-linux-generic64} feature同时启用\code{api-v0}并链接号表crate；\code{impl-dummy}不参与运行时ABI。

启用\code{impl-linux-generic64}后，聚合\code{syscall\_number}子模块导出\code{ActiveSyscallNumberTable = LinuxGeneric64}，供\code{wateros-syscall}与trap路径在编译期解析符号化调用号。

% 【本文件写什么】impl 实现层：dummy
%   - 定位：占位/链接桩实现，保证默认 feature 可编译。
%   - crate 路径：os/components/wateros-abi/abi-impl/impl-dummy/src/
%   - 如何实现对应 api-v0 trait：关键类型、状态机、数据结构
%   - 算法或硬件交互流程（可用伪代码/流程图）
%   - 本 impl 启用的 Cargo [features] 与 arch feature（impl-riscv64 等）
%   - 与其它 impl 变体的差异、切换 feature 时的影响
%   - 性能/内存假设与已知限制
% 【事实来源】
%   - 上述 crate 源码与 Cargo.toml
%   - os/feature-tree.txt 中指向本 impl 的 feature 链

\paragraph{职责}

\code{wateros-abi-impl-dummy}为Cargo工作区依赖解析与编译连通性占位，不包含真实平台ABI语义。对外仅提供\code{add(u64, u64) -> u64}模板函数，供crate内单测烟测使用。

\paragraph{与聚合层关系}

不通过\code{wateros-abi}聚合\code{lib.rs}导出；根\code{Cargo.toml}的\code{api-v0} feature链引用\code{impl-dummy/api-v0}仅为满足工作区成员解析。运行时主线选用\code{impl-linux-generic64}，切换至dummy不影响已启用号表的内核镜像。

\paragraph{限制}

无\code{SyscallNumberTable}实现，不参与trap或syscall分发。若仅启用\code{impl-dummy}而无\code{impl-linux-generic64}，\code{ActiveSyscallNumberTable}别名不存在，依赖活跃号表的crate无法编译。

% 【本文件写什么】impl 实现层：linux-generic64
% - 定位：Linux 兼容 generic64 ABI 实现，主线。
% - crate 路径：os/components/wateros-abi/abi-impl/impl-linux-generic64/src/
% - 如何实现对应 api-v0 trait：关键类型、状态机、数据结构
% - 算法或硬件交互流程，可用伪代码/流程图
% - 本 impl 启用的 Cargo [features] 与 arch feature，impl-riscv64 等
% - 与其它 impl 变体的差异、切换 feature 时的影响
% - 性能/内存假设与已知限制
% 【事实来源】
% - 上述 crate 源码与 Cargo.toml
% - os/feature-tree.txt 中指向本 impl 的 feature 链

\paragraph{LinuxGeneric64 号表}

零大小类型\code{LinuxGeneric64}实现\code{SyscallNumberTable}全部关联常量，数值来自Linux用户态可见的\code{asm-generic} 64位ABI。当前RISC-V64与LoongArch64共用此表，经根内核平台feature选择；若未来架构分叉须拆专用\code{impl-*}。

代表性映射，节选：\code{READ=63}、\code{WRITE=64}、\code{OPENAT=56}、\code{CLOSE=57}、\code{CLONE}经\code{FORK=220}映射、\code{EXEC=221}，即 \code{execve}、\code{EXIT=93}、\code{MMAP}族、\code{FUTEX}、\code{SOCKET}族、\code{EPOLL\_*}、\code{SYSLOG}等。xattr调用号5--16在表中单独列出，与\code{wateros-syscall}旁路分发一致。

\paragraph{特殊常量}

\code{SELECT}使用\code{usize::MAX}哨兵：\code{asm-generic} 64位无独立\code{select} nr，RISC-V64/LoongArch64用户态经\code{pselect6}，72实现\code{select}语义。\code{WAITPID=260}对应\code{wait4}；\code{FSTAT=80}覆盖\code{fstat}/\code{newfstatat} libc路径。

\paragraph{编译期校验}

crate内单测收集全部关联常量的\code{raw()}值，断言两两不等，防止号表碰撞。表项增多时须同步更新测试列表与\code{wateros-syscall}分发\code{match} arm。

\paragraph{与 syscall 边界}

\code{SyscallNumberTable}存在某号仅表示ABI编号可对齐；\code{wateros-syscall/impl-kernel}的\code{sys\_*}是否实现须单独核对。旁路号，含 \code{statx}、\code{fstatat}、\code{faccessat2} 等，在syscall层有额外\code{dispatch\_syscall\_aliases}分支，不在主号\code{match}重复。

\paragraph{切换与限制}

切换至其他号表impl须同时调整根内核feature链与所有使用\code{ActiveSyscallNumberTable}的\code{match}。当前无per-arch专用表，如riscv64-only；双架构共用是bring-up阶段的刻意简化。


